% !TEX program = pdflatex
% Draft v0.1 — venue TBD; using generic article layout, swap class later.
\documentclass[11pt]{article}
\usepackage[margin=1in]{geometry}
\usepackage{amsmath,amssymb}
\usepackage{graphicx}
\usepackage{booktabs}
\usepackage{multirow}
\usepackage{xcolor}
\usepackage[colorlinks=true,linkcolor=blue,citecolor=blue,urlcolor=blue]{hyperref}
\usepackage{CJKutf8} % for Chinese scenario names in examples; swap to xeCJK if compiling with xelatex

% ---- system name placeholder: final name TBD ----
\newcommand{\sysname}{\textsc{MemoirWorld}}
\newcommand{\todo}[1]{\textcolor{red}{[TODO: #1]}}

\title{\sysname: Consensus-Compressed Shared Memory for\\ Scalable Story-World Agent Societies}
\author{Yu Pan\thanks{\todo{author list, affiliation, email}}}
\date{}

\begin{document}
\maketitle

\begin{abstract}
Large language model (LLM) agent societies that role-play long-running story worlds
must remember what happened --- yet the dominant designs either store every
observation per agent, duplicating shared experiences across the cast, or
periodically summarize histories in ways that erase the very details that make
characters believable. We present \sysname, a multi-agent framework whose
long-term memory is a single shared vector store in which each atomic memory
entry carries an explicit \emph{owner set}. A normalization gate splits or
compresses every incoming memory into atomic, length-bounded statements, and a
\emph{consensus} mechanism merges semantically equivalent entries across agents,
keeping the shorter surface form, unioning owners, and preserving the earlier
story-time stamp. Combined with an event-driven tick scheduler in which idle
agents cost nothing, \sysname{} supports societies of dozens of characters,
including \emph{archived} agents for deceased characters whose memories remain
part of the shared record. A two-pass \emph{history sedimentation} pipeline
initializes a society directly from the first part of a novel (or a real-world
timeline): a registry pass canonicalizes characters, aliases, and places under
closed-world constraints, and a per-chapter, per-character exhaustive extraction
pass deposits time-stamped memories through the same consensus channel, so that
shared experiences (e.g., an oath sworn by three characters) converge to single
multi-owner entries. Agents then bootstrap their own goals from recalled
memories and simulate the \emph{sequel}, which we evaluate against the actual
continuation of the source. Across five scenarios --- three classical Chinese
novels, one English classic, and a real-world geopolitical timeline --- we
compare against Generative Agents, BookWorld, and G-Memory under identical
brains, budgets, and schedulers, measuring memory footprint, wall-clock and
token cost, narrative quality, continuation fidelity against ground truth, and
memory-grounded question answering. \sysname{} reaches comparable or better
quality at a fraction of the memory footprint and cost, and its memory dynamics
are directly interpretable: additions, merges, and revisions align with plot
events and character goals. \todo{fill in headline numbers}
\end{abstract}

\section{Introduction}
\label{sec:intro}

LLM-driven agent societies can now sustain multi-character role-play over long
horizons: characters converse, pursue goals, move through a world, and
accumulate experience. The bottleneck is memory. A society of $N$ characters
that each privately logs every turn stores the same shared scene $N$ times;
retrieval slows, context windows fill with near-duplicates, and cost grows with
the square of social activity. Conversely, aggressive summarization keeps
memory small but destroys the specific, character-bound details --- who swore
what to whom, and when --- that story worlds live on.

We argue that the unit of memory in a \emph{society} should not be the private
log but the \emph{shared fact}. Most of what characters remember, they remember
together: the oath in the peach garden is one fact with three owners, not three
private paraphrases. \sysname{} builds this observation into the storage layer
itself.

\paragraph{Contributions.}
\begin{enumerate}
  \item \textbf{Consensus-compressed shared memory.} A single vector store of
  atomic, length-bounded entries with explicit owner sets; inserts pass a
  normalization gate (split/compress into atoms) and a consensus check
  (embedding pre-filter + LLM equivalence) that merges duplicates across
  agents, keeps the shorter text, unions owners, and retains the earlier
  story-time. Deletion and revision are owner-set operations, so one agent
  forgetting does not erase a fact others still hold. (\S\ref{sec:memory})
  \item \textbf{History sedimentation.} A two-pass pipeline that turns the
  first part of a book or a real-world timeline into a society: a closed-world
  registry of canonical characters/places/aliases, then per-chapter,
  per-character exhaustive extraction of time-stamped atomic memories deposited
  through the consensus channel. Dead characters become \emph{archived} agents:
  never scheduled, but their memories persist in the shared record.
  Characters start the sequel with \emph{empty goal stacks} and bootstrap
  their own goals from recall --- goals are outcomes, not inputs.
  (\S\ref{sec:sedimentation})
  \item \textbf{An evaluation protocol for sequel simulation.} Beyond
  LLM-judged narrative quality, we score \emph{continuation fidelity} against
  the actual later part of the source (event hit rate, trajectory consistency)
  and \emph{memory-grounded QA} accuracy, and we report full memory/cost
  curves, enabling quality-per-byte comparisons. (\S\ref{sec:setup})
  \item \textbf{Evidence.} Across five diverse scenarios and three strong
  memory baselines under a strictly controlled runtime, \sysname{} attains
  similar or better quality with \todo{X$\times$} smaller memory and
  \todo{Y$\times$} lower cost, with interpretable memory dynamics.
  (\S\ref{sec:results}--\S\ref{sec:case})
\end{enumerate}

\section{Related Work}
\label{sec:related}

\paragraph{Memory for LLM agents.}
Generative Agents~\cite{park2023generative} introduced the memory stream:
append-only observations with recency/importance/relevance retrieval and
periodic reflection. MemGPT~\cite{packer2023memgpt} pages memory between
context and external storage. These designs are per-agent: in a society,
shared experiences are duplicated per participant. \todo{1-2 more recent
memory works (e.g., HippoRAG, MemoryBank) if venue expects}

\paragraph{Multi-agent societies and story worlds.}
BookWorld~\cite{chen2025bookworld} constructs role-agent societies from novels
with a world agent that periodically distills consensus summaries into world
memory while compressing per-role histories. G-Memory~\cite{gmemory2025}
maintains hierarchical interaction/insight structures shared across agents.
Our approach differs in \emph{where} deduplication happens: not as a
posterior summarization step over private logs, but at insert time in the
storage layer, preserving atomic facts with provenance (owner sets, story
time) instead of prose summaries. \todo{verify G-Memory citation details}

\paragraph{Simulation grounded in real timelines.}
\todo{brief para: social simulacra / world-model sims / wargaming with LLMs;
position the Russia--Ukraine scenario as continuation-fidelity benchmark
rather than forecasting claim.}

\section{The \sysname{} Framework}
\label{sec:framework}

\subsection{Runtime: agents, actions, and the tick barrier}
\label{sec:runtime}
An agent is one of three kinds --- \emph{character}, \emph{environment}, or
\emph{information carrier} --- with a pluggable brain (LLM, rules, or
retrieval). Each character holds a four-part short-term memory: a FIFO cache of
the last $K{=}20$ action--result pairs, a goal \emph{stack} (deeper goals more
fundamental), a status register with public/private keys, and a message inbox.
A global clock advances in ticks; in each tick every \emph{awake} agent decides
exactly one action (decisions run concurrently; effects apply in deterministic
agent order), messages sent at $t$ are delivered at $t{+}1$, and an agent whose
inbox and goal stack are both empty sleeps at zero cost until messaged.
Movement between environment nodes takes map-distance ticks. The action space
(21 actions) spans messaging (\texttt{say}, \texttt{gesture}), perception
(\texttt{observe}, \texttt{read}), cognition (\texttt{think},
\texttt{conclude}), goal and status management, memory operations
(\texttt{remember}, \texttt{recall}, \texttt{forget},
\texttt{revise\_memory}), and movement; agents are taught usage and standard
pipelines through a skill document injected into the system prompt.
\todo{figure: architecture diagram}

\subsection{Consensus-compressed shared long-term memory}
\label{sec:memory}
All long-term memory lives in one vector collection. An entry is a tuple
$(x, O, \tau, s)$: atomic text $x$ with $|x|\le L$ ($L{=}80$ chars by default),
owner set $O$, story-order $\tau$, and provenance $s$.

\paragraph{Normalization gate.} \texttt{remember}$(a, x)$ first checks $x$ for
length and atomicity (multiple sentence terminators or coordinating
connectives); violating inputs are split/compressed by one LLM call into atoms
of at most $L$ characters. Compliant inputs pass with zero LLM cost.

\paragraph{Consensus insert.} For each atom, retrieve the top-$k$ neighbors; if
none exceeds cosine similarity $\theta{=}0.86$, insert a new entry with
$O{=}\{a\}$. Otherwise a single LLM call judges semantic equivalence against
the candidates; on a match the store keeps the \emph{shorter} of the two
texts, sets $O \leftarrow O \cup \{a\}$, and keeps the \emph{earlier}
story-order (and its story-time label). Thus repeated tellings converge to one
canonical atom with provenance, and the shared-memory ratio
$|\{e : |O_e|\ge 2\}| / |E|$ becomes a measurable structural property of the
society.

\paragraph{Deletion and revision as owner-set operations.}
\texttt{forget} removes the caller from $O$ and deletes the entry only when
$O=\emptyset$; \texttt{revise\_memory} is forget-then-reinsert through the same
gate and consensus. \texttt{recall} is owner-filtered similarity search: agents
can only retrieve what they own, so information asymmetry between characters is
a storage-level invariant rather than a prompt-level convention.

\subsection{History sedimentation: initializing a society from a source text}
\label{sec:sedimentation}
Given the first part of a source (chapters $1..m$ of a novel; the first third
of a timeline), \sysname{} builds the society in two passes.
\textbf{Pass 1 (registry).} Chunk summaries feed one merge call that outputs a
closed-world registry: canonical IDs, display names, and \emph{aliases} for
characters, places, and information carriers. The registry is a reviewable
artifact; all later attribution is restricted to its IDs (unknown mentions are
either canonicalized through a resolver over names and aliases, added by a
single registry-completion call, or dropped with a warning --- never silently
invented).
\textbf{Pass 2 (sedimentation).} Chapter-aligned chunks are processed in story
order. For each chunk, a roster call lists the characters on stage and state
updates (location, alive); then \emph{one call per character} exhaustively
enumerates that character's experiences in the chunk --- actions, dialogue
gists, observations, relationship changes, gains and losses, stance shifts ---
as time-prefixed atoms; a coverage-audit call adds missed details. Every atom
is deposited via \texttt{remember} with story-order $\tau$, so shared scenes
merge into multi-owner entries at load time. Final states decide who is alive;
the dead become archived agents (never scheduled, not observable, memories
retained). Location handling is guarded by hard invariants (ASCII canonical
IDs, no duplicate names/aliases, every reference resolvable) enforced in code,
not prompts. The sediment is exported holographically (text + embeddings +
metadata), so societies rebuild with zero re-embedding.

\paragraph{Goal bootstrapping.} Sequel characters start with empty goal
stacks. On their first turns the view carries a bootstrap hint; agents recall
their past, observe their surroundings, conclude, and push their own
fundamental and immediate goals. In our runs, characters reconstructed
canonical motivations (e.g., \begin{CJK}{UTF8}{gbsn}曹操\end{CJK} deriving
``march on Xuzhou to avenge my father'' purely from sedimented memory and a
terse kickoff). \todo{tighten example after final runs}

\section{Experimental Setup}
\label{sec:setup}

\subsection{Scenarios}
\label{sec:scenarios}
Five societies spanning language, genre, cast size, and fiction vs.\ reality
(Table~\ref{tab:scenarios}). For each, the \emph{sediment} span initializes
memory; simulation then runs the sequel, and the held-out \emph{reference} span
provides ground truth for continuation fidelity. Sources are public-domain
texts (Wikisource) and CC-licensed timeline articles (Wikipedia), enabling
artifact release.

\begin{table}[t]
\centering\small
\begin{tabular}{lllll}
\toprule
Scenario & Language & Cast & Sediment & Reference \\
\midrule
Romance of the Three Kingdoms & zh & $\sim$40 & ch.~1--40 & ch.~41--60 \\
Dream of the Red Chamber & zh & $\sim$30 & ch.~1--40 & ch.~41--80 \\
Water Margin (scale stress) & zh & 108+ & ch.~1--40 & ch.~41--70 \\
War and Peace & en & $\sim$50 & Vol.~1--2 & Vol.~3 \\
Russia--Ukraine timeline & en & $\sim$25 & 2022/02--2023/06 & 2023/07--2024/12 \\
\bottomrule
\end{tabular}
\caption{Scenarios. Cast counts are registry sizes after curation; the
timeline scenario anthropomorphizes states, leaders, and organizations, with
key documents as information carriers. \todo{finalize counts}}
\label{tab:scenarios}
\end{table}

\subsection{Baselines}
\label{sec:baselines}
All baselines are re-implemented as \emph{memory backends} inside the \sysname{}
runtime: identical brains, action space, scheduler, kickoffs, and budgets ---
only the memory layer varies.
\textbf{Generative Agents}~\cite{park2023generative}: per-agent append-only
stream, recency/importance/relevance retrieval, periodic reflection.
\textbf{BookWorld}~\cite{chen2025bookworld}: per-role flat histories with
world-agent consensus summarization and post-hoc history compression.
\textbf{G-Memory}~\cite{gmemory2025}: shared hierarchical interaction/insight
memory. Sedimentation feeds each backend the same source span through its own
ingestion mechanism. \todo{document per-backend ingestion adaptations in appendix}

\subsection{Metrics}
\label{sec:metrics}
\textbf{Memory.} Entries and bytes over ticks (peak/final), per-agent vs.\
shared split, compression ratio against the raw event log.
\textbf{Time/cost.} Wall-clock per tick, retrieval latency, LLM calls and
tokens by purpose bucket; one-time sedimentation cost reported separately.
\textbf{Quality.}
(i) BookWorld's LLM-judge dimensions (anthropomorphism, character fidelity,
plot coherence, \todo{confirm exact dims}) for comparability;
(ii) \emph{Continuation fidelity}: an event checklist mined from the reference
span (e.g., ``Liu Bei relieves Tao Qian''), scored as hit rate by a judge with
human spot checks; blinded plausibility ratings of simulated vs.\ actual
continuations; per-character trajectory consistency;
(iii) \emph{Memory QA}: $N{=}100$ factual questions auto-generated from the
sediment span, answered only via \texttt{recall}; accuracy measures memory
faithfulness under compression.
Judges use a different model family than the simulated brains; all prompts and
rubrics released. \todo{judge model choice; inter-rater stats}

\subsection{Protocol}
Each scenario $\times$ backend: identical kickoff and tick budget (50 ticks),
3 seeds, mean$\pm$std. Primary figures are Pareto fronts of quality vs.\ final
memory size and vs.\ cumulative token cost. \todo{seeds/ticks may increase for
camera-ready}

\section{Results}
\label{sec:results}
\todo{Main table: 5 scenarios $\times$ 4 systems $\times$ (memory, cost, quality
triplet). Pareto figures. Scale-stress curve on Water Margin (memory vs.\ cast
size). Sleep-economy ablation showing cost vs.\ active-agent count.}

\section{Case Study: Interpretable Memory Dynamics}
\label{sec:case}
\todo{Three Kingdoms, 50-tick sequel: (a) timeline of one character's memory
adds/merges/revisions aligned against plot beats and the goal stack (Cao Cao:
revenge campaign); (b) a consensus merge example --- the peach-garden oath
converging to one three-owner entry; (c) memory heatmap (character $\times$
chapter) showing sediment density tracking narrative structure; (d) the
information-asymmetry anecdote: characters acting on stale beliefs until news
arrives, as an emergent property of owner-filtered recall.}

\section{Ablations}
\label{sec:ablations}
\todo{On Three Kingdoms (full grid) and consensus on/off everywhere:
consensus on/off; normalization gate on/off; $L \in \{40, 80, 160\}$;
recall top-$k \in \{3,5,10\}$; FIFO size $\{10,20,40\}$; goal bootstrap on/off;
archived agents kept vs.\ removed. Report memory/quality deltas.}

\section{Limitations and Ethics}
\label{sec:limitations}
Sedimentation quality depends on the extractor LLM; registry and life/death
states benefit from light human curation (we quantify curation effort).
The real-world scenario evaluates \emph{continuation fidelity on held-out past
events}; we make no forecasting claims, and we discuss the ethics of simulating
ongoing conflicts. \todo{expand}

\section{Conclusion}
\todo{write last}

\bibliographystyle{plain}
\begin{thebibliography}{9}
\bibitem{park2023generative} J.~S. Park et al. Generative Agents: Interactive
Simulacra of Human Behavior. UIST 2023.
\bibitem{packer2023memgpt} C.~Packer et al. MemGPT: Towards LLMs as Operating
Systems. arXiv:2310.08560.
\bibitem{chen2025bookworld} \todo{BookWorld: From Novels to Interactive Agent
Societies for Story Creation. ACL 2025. verify authors}
\bibitem{gmemory2025} \todo{G-Memory citation — verify title/venue/authors}
\end{thebibliography}

\end{document}
